\chapter{Метод}
\label{ch:metod}

\section{Формализиране на задачата}
\label{sec:formalizirane}

Нека една партия на ударни инструменти да се разгледа като редица от ноти.
За всяка нота $i$ разполагаме с
нейния \emph{структурно-времеви контекст} $\boldsymbol{x}_i$ — \emph{кой} каноничен глас свири,
кога (начален момент), в какъв стил и темпо, както и информация за съседните
удари. Търси се функция
\[
  f:\; \boldsymbol{x}_i \;\longmapsto\; v_i \in \{0,1,\dots,127\},
\]
която предсказва силата (\eng{velocity}) на нотата. В детерминистичната си постановка това
е задача за регресия; във вероятностната си постановка (раздел~\ref{sec:modeli})
целим цялото условно разпределение $p(v_i \mid \boldsymbol{x}_i)$, от което можем да
семплираме. Възможно е формулиране и на класификационна постановка,
вж. раздел~\ref{sec:prob-heads}.

\subsection{Строго условие: без изтичане на динамика}
\label{sec:no-leakage}

При извод входът е ,,плоско'' MIDI, чиито \eng{velocity} стойности \emph{ние
предсказваме}. Затова се налага абсолютно ограничение: \textbf{никоя сила
на коя да е нота не може да бъде вход при определяне силата на друга нота}. Това изключва цяла
категория характеристики (съседни динамики, текущи средни и др.) и целият
контекст $\boldsymbol{x}_i$ се изгражда само от структура и време. Пряко следствие е, че
последователният модел е \emph{неавторегресивен}: той вижда цялата структурна
редица и извежда \emph{всички} динамики наведнъж, в един паралелен пас — без
натрупване на грешка и без наредба за генериране.

\subsection{Защо без мрежа}
\label{sec:no-grid}

За разлика от мрежово-базираните груув модели (раздел~\ref{sec:svarzani-nevronni}),
тук \emph{не квантуваме} времето, а го подаваме като непрекъснат вход. Причините са две.
Мрежата е дуплетна и изкривява триолни/неравноделни усещания; а тъй
като времето при нас е условен вход, а не изход, изобщо няма нужда да го
дискретизираме. Вместо това всяка нота има своята \emph{непрекъсната метрична
фаза}, която представя точно всяка подразделба — триола при фаза прибл. $0{,}333$ или $0{,}667$,
квинтола при $0{,}2/0{,}4/\dots$, и суинг там, където реално е. Анализът в
раздел~\ref{sec:phase0} показва, че около $40\%$ от интервалите между удари
попадат \emph{измежду} шестнайсетичната мрежа — маса от орнаменти, триоли и суинг,
която мрежата би квантувала, сляла или изпуснала.

\section{Изследователска валидация на данните (Фаза 0)}
\label{sec:phase0}

Редица дефиниции, които иначе биха били произволни, са изведени по управляван от
данните начин чрез предварителен анализ \emph{само върху обучаващото} подмножество
($35\,217$ файла, $11{,}07$ млн.\ ноти, $0$ грешки при програмно прочитане на файловете).

\paragraph{Групиране на гласовете.} От хистограмите на използваните в набора данни
MIDI инструменти и сила (velocity) (фиг.~\ref{fig:pitch_counts},~\ref{fig:velocity_dists})
се извежда групиране в \textbf{14 канонични гласа}. Приложено е правилото
„слей редките и/или неразличимите, раздели честите и различимите'': томовете (темпаните), които са
неразличими по сила, се сливат в един каноничен глас; докато ударите по ръба на фуса
(\eng{hi-hat edge}) и електронният соло барабан (electric snare) — чести \emph{и}
разпределително различни (с десетки единици по-силни от аналозите си) — се пазят
като отделни гласове.

\paragraph{Толеранс за едновременност.} За да се определи кои удари са ,,практически
едновременни'' (един акорд), е нужно определяне на праг на интервалите между
началните моменти на нотите. Хистограмата на тези интервали (фиг.~\ref{fig:gap_hist})
показва две различни граници. В близост до нулата \emph{няма долина}: истинската едновременност
преминава монотонно в тесни орнаменти (форшлаг, тремоло), така че прагът не може да се
разчете като долина, а е \emph{документиран хиперпараметър}:
$\text{SIMULTANEITY\_TOL} = 0{,}02$ от текущото метрично време (около $11$~ms при медианното темпо),
избран в плоския спад далеч под най-малката музикална подразделба. Отделно, ясна
долина при $\approx0{,}149$ от метричното време отделя орнаментната маса от шестнайсетиновата
мрежа — това е доводът за представянето без мрежа (раздел~\ref{sec:no-grid}).

\begin{figure}[H]
  \centering
  \includegraphics[width=0.9\textwidth]{fig_gap_hist.png}
  \caption{Хистограма на интервалите между последователни начала (във времена). Липсва
  долина близо до нулата (едновременност $\to$ орнаменти), но има ясна долина при
  $\approx0{,}149$ от метричното време, отделяща под-шестнайсетиновата маса от мрежата
  ($0{,}25$ = шестнайсетина, $0{,}33$ = осмина триола, $0{,}5$ = осмина).}
  \label{fig:gap_hist}
\end{figure}

\section{Извличане на входни данни}
\label{sec:feature-extraction}

Всички характеристики са \emph{структурни}; никоя не използва \eng{velocity}
(вж. раздел~\ref{sec:no-leakage}). В настоящата работа се използват две свързани
представяния, стъпващи върху една и съща основа.

\subsection{Таблични характеристики (Раздел~A)}
\label{sec:features-a}

За табличния модел всяка нота е ред с ръчно проектиран контекстен прозорец. Тъй
като дърветата извършват имплицитен подбор, тук сме либерални:
\begin{itemize}
  \item \textbf{Глас} — каноничният глас, дефиниран във Фаза~0 (категорийна
        величина).
  \item \textbf{Метрична позиция} — $\sin$ и $\cos$ на фазата спрямо времето и такта,
        плюс самата фаза; отделно \emph{суинг \%} (изместване на слабото време) и
        \emph{най-близка подразделба} (осмина, шестнайсетина, 32-ина, осмина триола, и т.н.).
  \item \textbf{Общи за цялото изпълнение} — стил (жанр и пълен стилов низ), \eng{bpm},
        тактов размер, тип на такта.
  \item \textbf{Структурен контекст} — време до предишен/следващ удар (глобално и
        за същия глас) във времена, логаритмично мащабирано; едновременност (multi-hot кодиран
        вектор на гласовете, звучащи в един и същ момент, плюс техният брой); локална плътност в
        $\pm1$ време.
\end{itemize}

\subsection{Характеристики за невронните модели (Раздел~B)}
\label{sec:features-b}

За трансформъра се използва същата базова информация, представена нота по нота, но
\emph{ръчно проектираният прозорец се премахва} — модулът на ,,вниманието'' (\eng{self-attention})
сам научава контекста. Всяка нота е един \eng{token}: \eng{embedding} за
гласа ($\approx8$-мерен) и за жанра ($\approx16$-мерен), $\sin/\cos$ на метричната
фаза, стандартизирано (по статистики само от обучаващото подмножество) \eng{bpm} и
логаритмично мащабирано време до предишното събитие. Едновременните удари се
нареждат във входния вектор от ноти (токени) по височина, с флаг „същото начало като предишния''.
Тактовият размер се изпуска (Фаза~0: $98{,}78\%$ са $4/4$), тъй като е почти постоянен и до голяма
степен излишен спрямо фазата на такта.

\section{Модели}
\label{sec:modeli}

Върху един и същ конвейер от данни се изграждат и обучават три вида модели,
подредени по нарастваща сложност:

\subsection{Таблична база за сравнение — LightGBM}
\label{sec:lightgbm}

Като силен и интерпретируем базов модел се използва \eng{LightGBM} — градиентен
бустинг с дървета~\cite{ke2017lightgbm} — върху табличните признаци от
раздел~\ref{sec:features-a}. Моделът се обучава върху обучаващото подмножество с
ранно спиране по валидационното. Той би дал отговор на въпроса „заслужава ли
трансформърът своята сложност?'' за нашата задача.

\subsection{Последователен трансформър}
\label{sec:transformer}

Основният невронен модел е трансформър-енкодер~\cite{vaswani2017attention} върху
редицата от ноти, \emph{неавторегресивен} (раздел~\ref{sec:no-leakage}): той
предсказва всички динамики в един паралелен пас. Дълги партии се разделят на
прозорци с фиксирана дължина. Детерминистичната версия има линейна изходна глава,
връщаща една скаларна стойност, и се обучава с $L1$ функция на грешката (средна абсолютна
грешка).

\subsection{Вероятностни изходни глави}
\label{sec:prob-heads}

Точковата оценка на детерминистичния модел неизбежно „осреднява'' и заличава многомодалната вариация
(раздел~\ref{sec:dinamika}). Затова, \emph{върху същия обучен трансформърен гръбнак},
изходната глава се заменя така, че моделът да предсказва \emph{разпределение} над
\eng{velocity}, от което да семплираме при извод.
Три подвида модели от такъв тип, различаващи единствено по изходните глави,
биват обучени и сравнявани в настоящата работа:
\begin{itemize}
  \item \textbf{Гаусова глава} — предсказва нормално разпределение с параметри $(\mu,\sigma)$;
        обучава се с гаусова
        отрицателна логаритмична вероятност (\eng{NLL, Negative Log-Likelihood}).
  \item \textbf{Смес от гаусиани (MDN)}~\cite{bishop1994mdn} — $K$ тегла и $K$
        двойки $(\mu,\sigma)$; улавя многомодалността (акцент, призрачна нота, и т.н.).
  \item \textbf{Категорийна} — дискретни категории по изходната сила (\eng{velocity});
        обучава се с cross-entropy loss функция. Тъй като \eng{velocity} е ограничено
        цяло число, тази непараметрична глава не губи вероятностна маса извън
        $[0,127]$ и може да представи произволно многомодално разпределение
        (броят на категориите е хиперпараметър на системата).
\end{itemize}
При всяка от тези три глави можем да четем изхода \emph{детерминистично} (условно средно —
за точност и подредба) или чрез \emph{семплиране} (за реалистична човешка
вариация). Обучението използва \eng{AdamW}~\cite{kingma2015adam} като оптимизация.
